% Options for packages loaded elsewhere
\PassOptionsToPackage{unicode}{hyperref}
\PassOptionsToPackage{hyphens}{url}
\documentclass[
  ignorenonframetext,
]{beamer}
\newif\ifbibliography
\usepackage{pgfpages}
\setbeamertemplate{caption}[numbered]
\setbeamertemplate{caption label separator}{: }
\setbeamercolor{caption name}{fg=normal text.fg}
\beamertemplatenavigationsymbolsempty
% remove section numbering
\setbeamertemplate{part page}{
  \centering
  \begin{beamercolorbox}[sep=16pt,center]{part title}
    \usebeamerfont{part title}\insertpart\par
  \end{beamercolorbox}
}
\setbeamertemplate{section page}{
  \centering
  \begin{beamercolorbox}[sep=12pt,center]{section title}
    \usebeamerfont{section title}\insertsection\par
  \end{beamercolorbox}
}
\setbeamertemplate{subsection page}{
  \centering
  \begin{beamercolorbox}[sep=8pt,center]{subsection title}
    \usebeamerfont{subsection title}\insertsubsection\par
  \end{beamercolorbox}
}
% Prevent slide breaks in the middle of a paragraph
\widowpenalties 1 10000
\raggedbottom
\AtBeginPart{
  \frame{\partpage}
}
\AtBeginSection{
  \ifbibliography
  \else
    \frame{\sectionpage}
  \fi
}
\AtBeginSubsection{
  \frame{\subsectionpage}
}
\usepackage{iftex}
\ifPDFTeX
  \usepackage[T1]{fontenc}
  \usepackage[utf8]{inputenc}
  \usepackage{textcomp} % provide euro and other symbols
\else % if luatex or xetex
  \usepackage{unicode-math} % this also loads fontspec
  \defaultfontfeatures{Scale=MatchLowercase}
  \defaultfontfeatures[\rmfamily]{Ligatures=TeX,Scale=1}
\fi
\usepackage{lmodern}
\ifPDFTeX\else
  % xetex/luatex font selection
\fi
% Use upquote if available, for straight quotes in verbatim environments
\IfFileExists{upquote.sty}{\usepackage{upquote}}{}
\IfFileExists{microtype.sty}{% use microtype if available
  \usepackage[]{microtype}
  \UseMicrotypeSet[protrusion]{basicmath} % disable protrusion for tt fonts
}{}
\makeatletter
\@ifundefined{KOMAClassName}{% if non-KOMA class
  \IfFileExists{parskip.sty}{%
    \usepackage{parskip}
  }{% else
    \setlength{\parindent}{0pt}
    \setlength{\parskip}{6pt plus 2pt minus 1pt}}
}{% if KOMA class
  \KOMAoptions{parskip=half}}
\makeatother
\usepackage{color}
\usepackage{fancyvrb}
\newcommand{\VerbBar}{|}
\newcommand{\VERB}{\Verb[commandchars=\\\{\}]}
\DefineVerbatimEnvironment{Highlighting}{Verbatim}{commandchars=\\\{\}}
% Add ',fontsize=\small' for more characters per line
\newenvironment{Shaded}{}{}
\newcommand{\AlertTok}[1]{\textcolor[rgb]{1.00,0.00,0.00}{\textbf{#1}}}
\newcommand{\AnnotationTok}[1]{\textcolor[rgb]{0.38,0.63,0.69}{\textbf{\textit{#1}}}}
\newcommand{\AttributeTok}[1]{\textcolor[rgb]{0.49,0.56,0.16}{#1}}
\newcommand{\BaseNTok}[1]{\textcolor[rgb]{0.25,0.63,0.44}{#1}}
\newcommand{\BuiltInTok}[1]{\textcolor[rgb]{0.00,0.50,0.00}{#1}}
\newcommand{\CharTok}[1]{\textcolor[rgb]{0.25,0.44,0.63}{#1}}
\newcommand{\CommentTok}[1]{\textcolor[rgb]{0.38,0.63,0.69}{\textit{#1}}}
\newcommand{\CommentVarTok}[1]{\textcolor[rgb]{0.38,0.63,0.69}{\textbf{\textit{#1}}}}
\newcommand{\ConstantTok}[1]{\textcolor[rgb]{0.53,0.00,0.00}{#1}}
\newcommand{\ControlFlowTok}[1]{\textcolor[rgb]{0.00,0.44,0.13}{\textbf{#1}}}
\newcommand{\DataTypeTok}[1]{\textcolor[rgb]{0.56,0.13,0.00}{#1}}
\newcommand{\DecValTok}[1]{\textcolor[rgb]{0.25,0.63,0.44}{#1}}
\newcommand{\DocumentationTok}[1]{\textcolor[rgb]{0.73,0.13,0.13}{\textit{#1}}}
\newcommand{\ErrorTok}[1]{\textcolor[rgb]{1.00,0.00,0.00}{\textbf{#1}}}
\newcommand{\ExtensionTok}[1]{#1}
\newcommand{\FloatTok}[1]{\textcolor[rgb]{0.25,0.63,0.44}{#1}}
\newcommand{\FunctionTok}[1]{\textcolor[rgb]{0.02,0.16,0.49}{#1}}
\newcommand{\ImportTok}[1]{\textcolor[rgb]{0.00,0.50,0.00}{\textbf{#1}}}
\newcommand{\InformationTok}[1]{\textcolor[rgb]{0.38,0.63,0.69}{\textbf{\textit{#1}}}}
\newcommand{\KeywordTok}[1]{\textcolor[rgb]{0.00,0.44,0.13}{\textbf{#1}}}
\newcommand{\NormalTok}[1]{#1}
\newcommand{\OperatorTok}[1]{\textcolor[rgb]{0.40,0.40,0.40}{#1}}
\newcommand{\OtherTok}[1]{\textcolor[rgb]{0.00,0.44,0.13}{#1}}
\newcommand{\PreprocessorTok}[1]{\textcolor[rgb]{0.74,0.48,0.00}{#1}}
\newcommand{\RegionMarkerTok}[1]{#1}
\newcommand{\SpecialCharTok}[1]{\textcolor[rgb]{0.25,0.44,0.63}{#1}}
\newcommand{\SpecialStringTok}[1]{\textcolor[rgb]{0.73,0.40,0.53}{#1}}
\newcommand{\StringTok}[1]{\textcolor[rgb]{0.25,0.44,0.63}{#1}}
\newcommand{\VariableTok}[1]{\textcolor[rgb]{0.10,0.09,0.49}{#1}}
\newcommand{\VerbatimStringTok}[1]{\textcolor[rgb]{0.25,0.44,0.63}{#1}}
\newcommand{\WarningTok}[1]{\textcolor[rgb]{0.38,0.63,0.69}{\textbf{\textit{#1}}}}
\setlength{\emergencystretch}{3em} % prevent overfull lines
\providecommand{\tightlist}{%
  \setlength{\itemsep}{0pt}\setlength{\parskip}{0pt}}
\usepackage{bookmark}
\IfFileExists{xurl.sty}{\usepackage{xurl}}{} % add URL line breaks if available
\urlstyle{same}
\hypersetup{
  hidelinks,
  pdfcreator={LaTeX via pandoc}}

\author{\texorpdfstring{}{}}
\date{}

\begin{document}

\begin{frame}[fragile]{Reproducibility}
\protect\phantomsection\label{reproducibility}
Reproducible computational research requires version-pinned tools, fixed
inputs, and documented outputs {[}@peng2011reproducible{]}. COGANT
contributes the \textbf{graph-generation} slice of that story: identical
source trees and identical pipeline configuration should yield identical
exported bundles modulo declared nondeterminism (for example optional
neural embeddings if enabled).

\begin{block}{Version pinning}
\protect\phantomsection\label{version-pinning}
Pin:

\begin{itemize}
\tightlist
\item
  The \textbf{COGANT} version (\texttt{pyproject.toml} / package
  \texttt{\_\_version\_\_}) and Git commit when installing from source.
\item
  The \textbf{Python} minor version.
\item
  \textbf{Optional} Rust toolchain commit when building native
  extensions from \texttt{../cogant/rust/}.
\end{itemize}
\end{block}

\begin{block}{Artifact layout}
\protect\phantomsection\label{artifact-layout}
Pipeline output typically includes JSON for intermediate IRs,
\textbf{Generalized Notation Notation (GNN)} bundles (canonical
\texttt{model.gnn.md} plus companion JSON and interop artifacts),
validation reports, and optional HTML sites. Paths under the chosen
\texttt{output\_dir} should be treated as disposable but
\textbf{checksumable}: store hashes alongside published datasets derived
from COGANT exports when releasing research artifacts.
\end{block}

\begin{block}{Determinism}
\protect\phantomsection\label{determinism}
Parsing and graph construction aim for deterministic ordering on a fixed
filesystem snapshot. Features that pull in external models (for example
optional name or documentation embeddings consumed by the Generalized
Notation Notation exporter) introduce variability unless models and
seeds are fixed; the Generalized Notation Notation export document
(\texttt{../cogant/docs/GNN\_EXPORT.md}) calls out embedding dimensions
and optional behavior.
\end{block}

\begin{block}{Relation to the template repository}
\protect\phantomsection\label{relation-to-the-template-repository}
While this project remains outside
\href{../../../projects/}{\texttt{../../../projects/}}, it is
\textbf{not} executed by the root \texttt{./run.sh} discovery layer.
After promotion to
\href{../../../projects/cogant/}{\texttt{../../../projects/cogant/}}
with \texttt{src/}, \texttt{tests/}, and \texttt{pyproject.toml} per
template rules, the standard manuscript validation and PDF stages apply;
until then, validate Markdown from the template repository root,
e.g.~\texttt{uv\ run\ python\ -m\ infrastructure.validation.cli\ markdown\ ./projects\_in\_progress/cogant/manuscript/}.
\end{block}

\begin{block}{Validation gates}
\protect\phantomsection\label{validation-gates}
The pipeline enforces quality through three complementary validation
checkers, each targeting a different failure mode:

\textbf{IntegrityChecker.} Verifies structural soundness of the program
graph: all edge endpoints reference existing nodes, no unintended
duplicate nodes exist, orphaned nodes (zero in-degree and zero
out-degree) are flagged, and self-loops are reported unless explicitly
allowed by configuration. The checker also ensures that confidence
scores fall within \([0, 1]\) and that provenance records are non-empty
for every node and edge. A graph that fails integrity checks receives a
FAIL validation status; downstream export is blocked.

\textbf{SchemaChecker.} Validates each IR artifact against its JSON
schema (versioned alongside the COGANT package). Schema violations --
such as missing required fields, incorrect types, or unknown enum values
in \texttt{NodeKind} or \texttt{SemanticRole} -- are classified as ERROR
or FATAL depending on severity. Schema versions are recorded in the
validation report so that consumers can verify compatibility with their
import code.

\textbf{ProvenanceChecker.} Audits the provenance chain: every assertion
in the semantic mapping must trace back to at least one evidence source
(SourceCode, TypeSystem, ControlFlow, Heuristic, or External). The
checker flags mappings whose provenance is empty or whose confidence
score is inconsistent with the declared evidence tier -- for example, a
STATIC\_PLUS\_RUNTIME tier with no runtime trace evidence. These flags
appear as warnings rather than errors, since partial provenance is
expected for heuristic rules.

Together, these gates ensure that exported bundles meet a minimum
quality bar before reaching downstream models. The thresholds are
configurable (see \texttt{../cogant/docs/VALIDATION.md}); defaults
require \(\geq 95\%\) coverage, mean confidence \(\geq 0.85\), and zero
schema violations.
\end{block}

\begin{block}{Pipeline checkpoint and resume behavior}
\protect\phantomsection\label{pipeline-checkpoint-and-resume-behavior}
The pipeline writes a checkpoint file after each successfully completed
stage. The checkpoint records the stage name, completion timestamp,
input hash, and output hash. On a subsequent invocation with
\texttt{-\/-resume} (or \texttt{PipelineConfig.resume\ =\ True}), the
runner reads the checkpoint file, verifies that the input hash for each
completed stage still matches the current source tree, and skips stages
whose outputs are already valid. If a source file has changed, the
runner invalidates the affected stage and all downstream stages, then
re-executes from the earliest invalidated point.

Checkpoints are stored as JSON under the configured
\texttt{output\_dir}:

\begin{verbatim}
output/
  .cogant_checkpoint.json
\end{verbatim}

This mechanism is particularly useful during iterative development:
after an initial full run, changing a single source file triggers only
re-parsing, graph construction, and downstream stages for the affected
subgraph, rather than re-analyzing the entire repository.
\end{block}

\begin{block}{Output manifest structure}
\protect\phantomsection\label{output-manifest-structure}
Each pipeline run produces a manifest file (\texttt{manifest.json})
alongside the exported artifacts. The manifest provides a
machine-readable inventory of everything the run produced:

\begin{Shaded}
\begin{Highlighting}[]
\FunctionTok{\{}
  \DataTypeTok{"cogant\_version"}\FunctionTok{:} \StringTok{"0.1.0"}\FunctionTok{,}
  \DataTypeTok{"run\_id"}\FunctionTok{:} \StringTok{"run\_20260408\_143012"}\FunctionTok{,}
  \DataTypeTok{"timestamp"}\FunctionTok{:} \StringTok{"2026{-}04{-}08T14:30:12Z"}\FunctionTok{,}
  \DataTypeTok{"input\_hash"}\FunctionTok{:} \StringTok{"sha256:a1b2c3..."}\FunctionTok{,}
  \DataTypeTok{"config\_hash"}\FunctionTok{:} \StringTok{"sha256:d4e5f6..."}\FunctionTok{,}
  \DataTypeTok{"stages\_completed"}\FunctionTok{:} \OtherTok{[}\StringTok{"discovery"}\OtherTok{,} \StringTok{"parsing"}\OtherTok{,} \StringTok{"repo\_ir"}\OtherTok{,} \StringTok{"graph"}\OtherTok{,}
                        \StringTok{"translate"}\OtherTok{,} \StringTok{"validate"}\OtherTok{,} \StringTok{"export"}\OtherTok{]}\FunctionTok{,}
  \DataTypeTok{"stages\_skipped"}\FunctionTok{:} \OtherTok{[}\StringTok{"statespace"}\OtherTok{,} \StringTok{"process"}\OtherTok{]}\FunctionTok{,}
  \DataTypeTok{"artifacts"}\FunctionTok{:} \OtherTok{[}
    \FunctionTok{\{}\DataTypeTok{"path"}\FunctionTok{:} \StringTok{"model.gnn.json"}\FunctionTok{,} \DataTypeTok{"format"}\FunctionTok{:} \StringTok{"json"}\FunctionTok{,} \DataTypeTok{"size\_bytes"}\FunctionTok{:} \DecValTok{2048576}\FunctionTok{,}
     \DataTypeTok{"hash"}\FunctionTok{:} \StringTok{"sha256:..."}\FunctionTok{\}}\OtherTok{,}
    \FunctionTok{\{}\DataTypeTok{"path"}\FunctionTok{:} \StringTok{"model.pyg.pt"}\FunctionTok{,}   \DataTypeTok{"format"}\FunctionTok{:} \StringTok{"pytorch\_geometric"}\FunctionTok{,} \DataTypeTok{"size\_bytes"}\FunctionTok{:} \DecValTok{1024000}\FunctionTok{,}
     \DataTypeTok{"hash"}\FunctionTok{:} \StringTok{"sha256:..."}\FunctionTok{\}}\OtherTok{,}
    \FunctionTok{\{}\DataTypeTok{"path"}\FunctionTok{:} \StringTok{"validation\_report.json"}\FunctionTok{,} \DataTypeTok{"format"}\FunctionTok{:} \StringTok{"json"}\FunctionTok{,} \DataTypeTok{"size\_bytes"}\FunctionTok{:} \DecValTok{8192}\FunctionTok{,}
     \DataTypeTok{"hash"}\FunctionTok{:} \StringTok{"sha256:..."}\FunctionTok{\}}
  \OtherTok{]}\FunctionTok{,}
  \DataTypeTok{"statistics"}\FunctionTok{:} \FunctionTok{\{}
    \DataTypeTok{"nodes"}\FunctionTok{:} \DecValTok{320}\FunctionTok{,}
    \DataTypeTok{"edges"}\FunctionTok{:} \DecValTok{1250}\FunctionTok{,}
    \DataTypeTok{"mean\_confidence"}\FunctionTok{:} \FloatTok{0.92}\FunctionTok{,}
    \DataTypeTok{"validation\_status"}\FunctionTok{:} \StringTok{"PASS"}
  \FunctionTok{\}}
\FunctionTok{\}}
\end{Highlighting}
\end{Shaded}

Storing per-artifact SHA-256 hashes enables consumers to verify that
exported bundles have not been modified after generation. When
publishing datasets derived from COGANT exports, including the manifest
alongside the data satisfies the reproducibility requirements outlined
by {[}@peng2011reproducible{]}: any researcher with the same source
tree, COGANT version, and configuration can regenerate the artifacts and
compare hashes.
\end{block}

\begin{block}{Data ethics and licensing}
\protect\phantomsection\label{data-ethics-and-licensing}
Exported graphs can contain identifiers and comments from source code.
Redistribution of derived graphs must respect the licenses of input
repositories and organizational data policies.
\end{block}
\end{frame}

\end{document}
